\documentclass[11pt]{article}

\usepackage[margin=1in]{geometry}
\usepackage{amsmath,amssymb,amsthm}
\usepackage{mathtools}
\usepackage{microtype}
\usepackage[hidelinks]{hyperref}

\newtheorem{theorem}{Theorem}
\newtheorem{lemma}{Lemma}

\newcommand{\R}{\mathbb R}
\newcommand{\Pp}{\mathsf P}
\newcommand{\Qq}{\mathsf Q}
\newcommand{\E}{\mathbb E}

\title{Approximate Differential Privacy Is Closed Under Joint Weak Limits}
\author{}
\date{}

\begin{document}
\maketitle

\section{Statement}

Let \(Y\) be a second-countable metric space equipped with its Borel
sigma-algebra.  For probability measures \(\Pp,\Qq\) on \(Y\), define the
ordered \((\varepsilon,\delta)\)-privacy relation by
\begin{equation}
  \Pp(A)
  \leq
  e^\varepsilon \Qq(A)+\delta
  \qquad
  \text{for every Borel set }A\subseteq Y.
  \label{eq:ordered-dp}
\end{equation}
This is one ordered direction of approximate differential privacy.  The
reverse direction is obtained by interchanging \(\Pp\) and \(\Qq\).

\begin{theorem}[Ordered weak-limit closure]
\label{thm:ordered-closure}
Let \((\Pp_i)_i\) and \((\Qq_i)_i\) be nets of probability measures on \(Y\),
indexed along a nontrivial filter.  Assume
\[
  \Pp_i \Rightarrow \Pp,
  \qquad
  \Qq_i \Rightarrow \Qq,
\]
and suppose that, eventually in \(i\),
\[
  \Pp_i(A)\leq e^\varepsilon\Qq_i(A)+\delta
  \qquad
  \text{for every Borel }A.
\]
If \(\delta\geq0\), then
\[
  \Pp(A)\leq e^\varepsilon\Qq(A)+\delta
  \qquad
  \text{for every Borel }A.
\]
\end{theorem}

The use of nets rather than sequences is only a topological generalization.
For the DP-Adam application the index is \(\mathbb N\) and the filter is
eventual convergence as the mesh is refined.

\section{Proof}

The direct Portmanteau comparison of one set under both measures is not
enough.  For a closed set \(F\), weak convergence gives an upper bound for
\(\Pp(F)\), while for the same closed set it gives an upper bound, rather than
the required lower control, for \(\Qq(F)\).  We therefore pass first through
bounded continuous test functions.

\begin{lemma}[From events to bounded tests]
\label{lem:event-to-test}
Suppose \(\Pp\) and \(\Qq\) satisfy~\eqref{eq:ordered-dp}.  If
\(f:Y\to[0,1]\) is measurable, then
\[
  \int f\,d\Pp
  \leq
  e^\varepsilon\int f\,d\Qq+\delta.
\]
\end{lemma}

\begin{proof}
Use the layer-cake representation
\[
  \int f\,d\Pp
  =
  \int_0^1 \Pp\{f>t\}\,dt.
\]
Apply~\eqref{eq:ordered-dp} to every level set ({f>t}), then integrate:
\[
  \int_0^1 \Pp\{f>t\}\,dt
  \leq
  e^\varepsilon\int_0^1 \Qq\{f>t\}\,dt
  +\int_0^1\delta\,dt.
\]
The two layer-cake identities and the fact that the interval has length one
give the claim.
\end{proof}

Now let \(f:Y\to[0,1]\) be bounded and continuous.  By
Lemma~\ref{lem:event-to-test}, eventually
\[
  \int f\,d\Pp_i
  \leq
  e^\varepsilon\int f\,d\Qq_i+\delta.
\]
Weak convergence is exactly convergence of integrals against bounded
continuous functions.  Passing to the limit yields
\begin{equation}
  \int f\,d\Pp
  \leq
  e^\varepsilon\int f\,d\Qq+\delta
  \qquad (0\leq f\leq1,\ f\text{ continuous}).
  \label{eq:continuous-test}
\end{equation}

It remains to recover arbitrary Borel events.  Let \(F\subseteq U\), where
\(F\) is closed and \(U\) is open.  Metric spaces are normal, so there is a
continuous function \(f:Y\to[0,1]\) satisfying
\[
  f=1\text{ on }F,
  \qquad
  f=0\text{ on }U^c.
\]
Consequently,
\[
  \Pp(F)
  \leq \int f\,d\Pp
  \leq e^\varepsilon\int f\,d\Qq+\delta
  \leq e^\varepsilon\Qq(U)+\delta.
\]

Finally let \(A\) be Borel.  Borel probability measures on a
second-countable metric space are regular.  Thus
\[
  \Pp(A)=\sup_{F\subseteq A,\,F\text{ closed}}\Pp(F),
  \qquad
  \Qq(A)=\inf_{A\subseteq U,\,U\text{ open}}\Qq(U).
\]
Apply the preceding closed--open inequality to every
\(F\subseteq A\subseteq U\), then take the supremum over \(F\) and the
infimum over \(U\).  This proves Theorem~\ref{thm:ordered-closure}.

\section{Mechanism corollary}

Let \(\mathcal D\) be a dataset type with an ordered adjacency relation
\(\sim\).  For every index \(i\), let
\[
  M_i:\mathcal D\to\mathcal P(Y)
\]
be a probability-measure mechanism, and let
\(M_\infty:\mathcal D\to\mathcal P(Y)\).  Assume:
\begin{enumerate}
  \item for every fixed dataset \(D\), \(M_i(D)\Rightarrow M_\infty(D)\);
  \item eventually in \(i\), the complete mechanism \(M_i\) is
  \((\varepsilon,\delta)\)-DP, with the same index tail and the same privacy
  parameters for all ordered adjacent pairs.
\end{enumerate}
Then \(M_\infty\) is \((\varepsilon,\delta)\)-DP.  Indeed, fix an ordered
adjacent pair \(D\sim D'\) and apply Theorem~\ref{thm:ordered-closure} to
\(M_i(D)\) and \(M_i(D')\).  Since the definition quantifies over ordered
pairs, no unstated symmetry argument is needed.

The uniform eventual-mechanism hypothesis is stronger than saying that every
adjacent pair has its own unrelated tail.  It is exactly what is supplied by
the finite-mechanism accountant in the DP-Adam application.

\section{Formal source crosswalk and scope}

The formal development is independently replayed in the repository.  The
main declarations are:
\begin{center}
\begin{tabular}{p{0.43\linewidth}p{0.47\linewidth}}
\hline
Mathematical role & Lean declaration \\
\hline
Ordered event inequality &
\path{Mf.DP.OrderedApproxDPLawPair} \\
Event inequality for bounded tests &
\path{OrderedApproxDPLawPair.lintegral_le} \\
Passage of continuous tests to the limit &
\path{orderedApproxDPLawPair_limit_lintegral_le} \\
Closed--open separation &
\path{orderedApproxDPLawPair_closed_le_open_of_limit_lintegral} \\
Recovery of Borel events &
\path{orderedApproxDPLawPair_of_limit_lintegral} \\
Ordered closure theorem &
\path{orderedApproxDPLawPair_closed_under_joint_weak_convergence} \\
Mechanism-level corollary &
\path{approxDP_closed_under_pointwise_weak_convergence} \\
\hline
\end{tabular}
\end{center}

The theorem is qualitative: it gives no finite-index privacy correction and
no rate of weak convergence.  It also assumes probability laws on a common
second-countable metric Borel carrier.  It does not identify a mathematical
mechanism with any software implementation.  Those are separate application
obligations.

The printed axiom boundary of the terminal Lean theorems contains only
\[
  \texttt{propext},
  \qquad
  \texttt{Classical.choice},
  \qquad
  \texttt{Quot.sound}.
\]

\end{document}
